% Source: Weinberg GR, physical PDF pp. 125--128 (printed pp. 100--103).
% Coverage: Section 4.5, Transformation of the Affine Connection; equations
%           (4.5.1)--(4.5.11).
% Figures/tables/footnotes: none.
% Status: source-reviewed and compile-clean.
% Uncertainties: none.

\subsection{Transformation of the Affine Connection}\label{sec:4.5}

Apart from the rather trivial example of the tensor densities, there appears
throughout the laws of physics one other very important nontensor, the affine
connection. We recall its definition,
\begin{equation}
  \Gamma^\lambda{}_{\mu\nu}
  =
  \frac{\partial x^\lambda}{\partial \xi^\alpha}
  \frac{\partial^2\xi^\alpha}
       {\partial x^\mu\,\partial x^\nu},
  \tag{4.5.1}\label{eq:4.5.1}
\end{equation}
where \(\xi^\alpha(x)\) is a locally inertial coordinate system. Passing from
\(x^\mu\) to a different system \(x'^\mu\), we find that
\begin{equation}
  \Gamma'^\lambda{}_{\mu\nu}
  =
  \frac{\partial x'^\lambda}{\partial x^\rho}
  \frac{\partial x^\tau}{\partial x'^\mu}
  \frac{\partial x^\sigma}{\partial x'^\nu}
  \Gamma^\rho{}_{\tau\sigma}
  +
  \frac{\partial x'^\lambda}{\partial x^\rho}
  \frac{\partial^2x^\rho}
       {\partial x'^\mu\,\partial x'^\nu}.
  \tag{4.5.2}\label{eq:4.5.2}
\end{equation}
The first term on the right is what we would expect if
\(\Gamma^\lambda{}_{\mu\nu}\) were a tensor; the second term is
inhomogeneous, and makes it a nontensor.

Tensor analysis provides a very simple way of establishing the relation
between \(\Gamma^\lambda{}_{\mu\nu}\) and \(g_{\mu\nu}\). Differentiating
the transformation law \eqref{eq:4.2.6} for the metric and taking the
combination
\(\partial'_\mu g'_{\kappa\nu}+\partial'_\nu g'_{\kappa\mu}
-\partial'_\kappa g'_{\mu\nu}\), we find
\begin{equation}
  \Gamma[g]'^\lambda{}_{\mu\nu}
  =
  \frac{\partial x'^\lambda}{\partial x^\rho}
  \frac{\partial x^\tau}{\partial x'^\mu}
  \frac{\partial x^\sigma}{\partial x'^\nu}
  \Gamma[g]^\rho{}_{\tau\sigma}
  +
  \frac{\partial x'^\lambda}{\partial x^\rho}
  \frac{\partial^2x^\rho}
       {\partial x'^\mu\,\partial x'^\nu},
  \tag{4.5.3}\label{eq:4.5.3}
\end{equation}
where the modern notation \(\Gamma[g]\) denotes the connection built from the
metric,
\begin{equation}
  \Gamma[g]^\lambda{}_{\mu\nu}
  \equiv
  \frac12 g^{\lambda\kappa}
  \left(
    \partial_\mu g_{\kappa\nu}
    +\partial_\nu g_{\kappa\mu}
    -\partial_\kappa g_{\mu\nu}
  \right).
  \tag{4.5.4}\label{eq:4.5.4}
\end{equation}
Subtracting Eq.~\eqref{eq:4.5.3} from Eq.~\eqref{eq:4.5.2}, we see that
\(\Gamma^\lambda{}_{\mu\nu}-\Gamma[g]^\lambda{}_{\mu\nu}\) is a tensor:
\begin{equation}
  \left[
    \Gamma^\lambda{}_{\mu\nu}
    -\Gamma[g]^\lambda{}_{\mu\nu}
  \right]'
  =
  \frac{\partial x'^\lambda}{\partial x^\rho}
  \frac{\partial x^\tau}{\partial x'^\mu}
  \frac{\partial x^\sigma}{\partial x'^\nu}
  \left[
    \Gamma^\rho{}_{\tau\sigma}
    -\Gamma[g]^\rho{}_{\tau\sigma}
  \right].
  \tag{4.5.5}\label{eq:4.5.5}
\end{equation}

The equivalence principle tells us that there is a special coordinate system
\(\xi^\alpha\) in which, at a given point \(X\), the effects of gravitation
are absent. In this system there can be no gravitational force on free
particles, so \(\Gamma^\lambda{}_{\mu\nu}\) vanishes, and there can be no
gravitational red shift between infinitesimally separated points, so the
first derivatives of \(g_{\mu\nu}\) vanish. Since
\(\Gamma^\rho{}_{\tau\sigma}-\Gamma[g]^\rho{}_{\tau\sigma}\) vanishes in a
locally inertial coordinate system, and since it is a tensor, it must vanish
in all coordinate systems; that is,
\begin{equation}
  \Gamma^\lambda{}_{\mu\nu}
  =
  \Gamma[g]^\lambda{}_{\mu\nu}.
  \tag{4.5.6}\label{eq:4.5.6}
\end{equation}
In the following we therefore write simply
\(\Gamma^\lambda{}_{\mu\nu}\) for the Levi-Civita connection
\eqref{eq:4.5.4}.

It is useful to have at hand an alternative formula for the inhomogeneous
term in the transformation rule of \(\Gamma^\lambda{}_{\mu\nu}\).
Differentiate the identity
\[
  \frac{\partial x'^\lambda}{\partial x^\rho}
  \frac{\partial x^\rho}{\partial x'^\nu}
  =\delta^\lambda{}_\nu
\]
with respect to \(x'^\mu\); we find immediately that
\begin{equation}
  \frac{\partial x'^\lambda}{\partial x^\rho}
  \frac{\partial^2x^\rho}
       {\partial x'^\mu\,\partial x'^\nu}
  =
  -
  \frac{\partial x^\rho}{\partial x'^\nu}
  \frac{\partial x^\sigma}{\partial x'^\mu}
  \frac{\partial^2x'^\lambda}
       {\partial x^\rho\,\partial x^\sigma}.
  \tag{4.5.7}\label{eq:4.5.7}
\end{equation}
We can therefore write Eq.~\eqref{eq:4.5.2} as
\begin{equation}
  \Gamma'^\lambda{}_{\mu\nu}
  =
  \frac{\partial x'^\lambda}{\partial x^\rho}
  \frac{\partial x^\tau}{\partial x'^\mu}
  \frac{\partial x^\sigma}{\partial x'^\nu}
  \Gamma^\rho{}_{\tau\sigma}
  -
  \frac{\partial x^\rho}{\partial x'^\nu}
  \frac{\partial x^\sigma}{\partial x'^\mu}
  \frac{\partial^2x'^\lambda}
       {\partial x^\rho\,\partial x^\sigma}.
  \tag{4.5.8}\label{eq:4.5.8}
\end{equation}
This is just what we would have found by first performing the inverse
transformation \(x'\to x\), and then solving for
\(\Gamma'^\lambda{}_{\mu\nu}\).

We are now in a position to use the Principle of General Covariance to give
an alternative proof that a freely falling particle obeys the equation of
motion
\begin{equation}
  \Dworldline{u^\mu}
  \equiv
  \frac{\dd u^\mu}{\dd\tau}
  +\Gamma^\mu{}_{\nu\lambda}u^\nu u^\lambda
  =0,
  \qquad
  u^\mu\equiv\frac{\dd x^\mu}{\dd\tau},
  \tag{4.5.9}\label{eq:4.5.9}
\end{equation}
where
\begin{equation}
  \dd\tau^2=-g_{\mu\nu}\,\dd x^\mu\dd x^\nu.
  \tag{4.5.10}\label{eq:4.5.10}
\end{equation}
First, note that Eqs.~\eqref{eq:4.5.9} and \eqref{eq:4.5.10} are true in
the absence of gravitation, because setting
\(\Gamma^\mu{}_{\nu\lambda}=0\) and \(g_{\mu\nu}=\eta_{\mu\nu}\) gives
\[
  \frac{\dd^2x^\mu}{\dd\tau^2}=0,
  \qquad
  \dd\tau^2=-\eta_{\mu\nu}\,\dd x^\mu\dd x^\nu,
\]
and these are the correct equations for a free particle in special
relativity.

Second, Eq.~\eqref{eq:4.5.10} is invariant under a general coordinate
transformation, while differentiating
\[
  u'^\mu
  =
  \frac{\partial x'^\mu}{\partial x^\nu}u^\nu
\]
and using Eq.~\eqref{eq:4.5.8} gives
\begin{equation}
  \Dworldline{u'^\mu}
  =
  \frac{\partial x'^\mu}{\partial x^\kappa}
  \Dworldline{u^\kappa}.
  \tag{4.5.11}\label{eq:4.5.11}
\end{equation}
Thus the left-hand side of Eq.~\eqref{eq:4.5.9} is a vector, and
Eqs.~\eqref{eq:4.5.9} and \eqref{eq:4.5.10} are manifestly covariant. The
Principle of General Covariance then tells us that they are true in general
gravitational fields, because, to repeat the reasoning of Section~\ref{sec:4.1},
they are true in all coordinate systems if true in any one system, and they
are true in the locally inertial coordinate systems.
